\documentclass[11pt,a4paper]{article}

\usepackage[utf8]{inputenc}
\usepackage[T1]{fontenc}
\usepackage[french]{babel}
\usepackage[a4paper,margin=1.8cm]{geometry}
\usepackage{amsmath,amssymb}
\usepackage{siunitx}
\usepackage{tabularx,array,booktabs,longtable,multirow}
\usepackage{enumitem}
\usepackage{xcolor}
\usepackage{fancyhdr}
\usepackage{lastpage}
\usepackage{hyperref}
\usepackage{tcolorbox}
\usepackage{titlesec}
\usepackage{multicol}
\usepackage{chemfig}
\usepackage{tikz}

\hypersetup{colorlinks=true,linkcolor=black,urlcolor=blue,citecolor=black}
\pagestyle{fancy}
\fancyhf{}
\lhead{Première générale -- Spécialité Physique-Chimie}
\rhead{Sujet type bac -- version longue}
\cfoot{\thepage/\pageref{LastPage}}

\titleformat{\section}{\large\bfseries}{}{0em}{}
\titleformat{\subsection}{\normalsize\bfseries}{}{0em}{}

\newcommand{\pts}[1]{\hfill{\small\textbf{(#1 pts)}}}
\newcommand{\bareme}[1]{\hfill{\small\textit{#1}}}
\setlist[enumerate]{itemsep=0.35em}

\begin{document}



\vspace{0.6em}

\begin{tcolorbox}[colback=blue!4,colframe=black,title=Conseils de méthode]
Dans les questions de spectroscopie IR, il faut toujours croiser les indices : \textbf{présence ou absence d'une bande}, \textbf{position du nombre d'onde}, \textbf{intensité et largeur éventuelles}, puis vérifier la cohérence avec la \textbf{famille chimique} et avec la \textbf{formule brute}. En cohésion de la matière, il faut raisonner avec les \textbf{liaisons polarisées}, la \textbf{géométrie de la molécule} et la nature des \textbf{interactions entre entités}. 
\end{tcolorbox}

\section*{Données utiles}

\begin{center}
\begin{tabular}{>{\raggedright\arraybackslash}p{4.8cm} >{\raggedright\arraybackslash}p{3.8cm} >{\raggedright\arraybackslash}p{6.2cm}}
\toprule
\textbf{Liaison ou groupe} & \textbf{Nombre d'onde} (cm$^{-1}$) & \textbf{Indication} \\
\midrule
O--H (alcool) & 3200 -- 3700 & bande large \\
O--H (acide carboxylique) & 2500 -- 3200 & très large, souvent étalée \\
C=O & 1650 -- 1750 & bande intense \\
C--H & 2850 -- 3100 & souvent présente \\
N--H & 3100 -- 3500 & bande souvent assez fine \\
\bottomrule
\end{tabular}
\end{center}

\begin{center}
\begin{tabular}{>{\raggedright\arraybackslash}p{5.7cm} >{\raggedright\arraybackslash}p{8.8cm}}
\toprule
\textbf{Famille} & \textbf{Suffixe principal} \\
\midrule
Alcool & -ol \\
Aldéhyde & -al \\
Cétone & -one \\
Acide carboxylique & acide ... -oïque \\
Ester & ... -oate de ... \\
\bottomrule
\end{tabular}
\end{center}

\begin{tcolorbox}[colback=gray!7,colframe=black,title=Rappels de cours utiles]
Une liaison covalente est \textbf{polarisée} si les électrons liants sont davantage attirés par un des deux atomes. Une molécule peut contenir des liaisons polarisées mais être \textbf{apolaire} si sa géométrie fait se compenser les effets de ces polarisations. Les liaisons hydrogène peuvent apparaître lorsqu'un atome d'hydrogène est lié à un atome très électronégatif comme l'oxygène et interagit avec un doublet d'un autre atome électronégatif.
\end{tcolorbox}

\newpage
\section*{Exercice 1 -- Reconnaître une transformation d'un alcool par IR \pts{6}}

On réalise au laboratoire la transformation d'une espèce organique A en une espèce organique B. On sait que A est un alcool et que la transformation correspond à une oxydation douce. Deux produits sont envisagés pour B :

\begin{itemize}[leftmargin=1.2em]
    \item le \textbf{propanal} ;
    \item la \textbf{propanone}.
\end{itemize}

Le spectre IR de B présente :
\begin{itemize}[leftmargin=1.2em]
    \item une bande intense vers \SI{1715}{cm^{-1}} ;
    \item aucune large bande entre \SI{3200}{cm^{-1}} et \SI{3700}{cm^{-1}}.
\end{itemize}

\begin{enumerate}[label=\textbf{\arabic*.},leftmargin=1.2em]
    \item Que signifie l'absence de large bande entre \SI{3200}{cm^{-1}} et \SI{3700}{cm^{-1}} ? \pts{1}
    \item À quel groupe caractéristique associe-t-on la bande vers \SI{1715}{cm^{-1}} ? \pts{1}
    \item Pourquoi peut-on affirmer que B n'est plus un alcool ? \pts{1}
    \item Le produit B est-il nécessairement un aldéhyde ? Justifier. \pts{1}
    \item Expliquer pourquoi les deux espèces proposées restent compatibles avec le spectre. \pts{1}
    \item Quelle information supplémentaire simple (autre que le spectre donné) permettrait de distinguer sûrement propanal et propanone ? \pts{1}
\end{enumerate}

\section*{Exercice 2 -- Cohésion de la matière : polarité et géométrie \pts{6}}

On s'intéresse à trois espèces moléculaires : l'eau H$_2$O, le dioxyde de carbone CO$_2$ et le sulfure d'hydrogène H$_2$S.

On admet les informations suivantes :
\begin{itemize}[leftmargin=1.2em]
    \item les liaisons O--H, C=O et S--H sont polarisées ;
    \item H$_2$O et H$_2$S sont des molécules coudées ;
    \item CO$_2$ est une molécule linéaire.
\end{itemize}

\begin{enumerate}[label=\textbf{\arabic*.},leftmargin=1.2em]
    \item Rappeler ce qu'est une liaison polarisée. \pts{1}
    \item Expliquer pourquoi H$_2$O est une molécule polaire. \pts{1.5}
    \item Expliquer pourquoi CO$_2$ peut être apolaire alors qu'il contient des liaisons polarisées. \pts{1.5}
    \item H$_2$S est-elle plutôt apolaire ou polaire ? Justifier. \pts{1}
    \item Parmi H$_2$O, CO$_2$ et H$_2$S, quelle espèce interagit le plus fortement avec l'eau liquide ? Justifier qualitativement. \pts{1}
\end{enumerate}

\section*{Exercice 3 -- Nomenclature et familles organiques \pts{5}}

On considère les espèces suivantes :
\begin{center}
\begin{tabular}{ccl}
A & : & CH$_3$--CH$_2$--OH \\
B & : & CH$_3$--CH$_2$--CHO \\
C & : & CH$_3$--CO--CH$_3$ \\
D & : & CH$_3$--CH$_2$--COOH \\
E & : & CH$_3$--COO--CH$_2$--CH$_3$ \\
\end{tabular}
\end{center}

\begin{enumerate}[label=\textbf{\arabic*.},leftmargin=1.2em]
    \item Donner le nom de chacune de ces espèces. \pts{2.5}
    \item Pour chacune, préciser la famille chimique. \pts{1.5}
    \item Parmi ces espèces, lesquelles possèdent un groupe carbonyle ? \pts{1}
\end{enumerate}

\section*{Exercice 4 -- Identifier une molécule à partir de sa formule brute et d'un spectre IR \pts{7}}

On cherche à identifier une espèce organique X de formule brute C$_3$H$_6$O.

Son spectre infrarouge présente :
\begin{itemize}[leftmargin=1.2em]
    \item une bande intense vers \SI{1720}{cm^{-1}} ;
    \item aucune large bande vers \SI{3300}{cm^{-1}}.
\end{itemize}

Trois espèces sont proposées :
\begin{itemize}[leftmargin=1.2em]
    \item propan-1-ol ;
    \item propanal ;
    \item propanone.
\end{itemize}

\begin{enumerate}[label=\textbf{\arabic*.},leftmargin=1.2em]
    \item Quel groupe caractéristique met en évidence la bande à \SI{1720}{cm^{-1}} ? \pts{1}
    \item Quelle famille peut-on éliminer immédiatement grâce à l'absence de large bande vers \SI{3300}{cm^{-1}} ? \pts{1}
    \item Expliquer pourquoi la formule brute C$_3$H$_6$O est compatible avec un aldéhyde ou une cétone. \pts{1}
    \item Identifier les deux espèces encore possibles. \pts{1}
    \item Le spectre IR suffit-il ici à départager ces deux espèces ? Expliquer. \pts{1}
    \item Donner un argument de nomenclature permettant de reconnaître la cétone parmi les noms proposés. \pts{1}
    \item Conclure : que peut-on affirmer avec certitude, et que ne peut-on pas affirmer ? \pts{1}
\end{enumerate}

\section*{Exercice 5 -- Solubilité, interactions et raisonnement comparatif \pts{8}}

On souhaite comparer le comportement dans l'eau de quatre espèces :
\begin{center}
\begin{tabular}{ccl}
M$_1$ & : & butane \\
M$_2$ & : & butan-1-ol \\
M$_3$ & : & propanone \\
M$_4$ & : & acide éthanoïque \\
\end{tabular}
\end{center}

On rappelle que l'eau est un solvant polaire et capable de former des liaisons hydrogène.

\begin{enumerate}[label=\textbf{\arabic*.},leftmargin=1.2em]
    \item Parmi ces quatre espèces, laquelle paraît la moins soluble dans l'eau ? Justifier. \pts{1.5}
    \item Quelle(s) espèce(s) peuvent former des liaisons hydrogène avec l'eau ? Justifier à partir de leur structure ou de leur groupe caractéristique. \pts{2}
    \item Expliquer pourquoi la présence d'un groupe O--H tend en général à favoriser la solubilité dans l'eau. \pts{1.5}
    \item La propanone ne possède pas de liaison O--H. Peut-elle quand même interagir avec l'eau ? Expliquer. \pts{1.5}
    \item Entre butan-1-ol et acide éthanoïque, laquelle des deux espèces semble interagir le plus fortement avec l'eau ? Proposer un raisonnement qualitatif. \pts{1.5}
\end{enumerate}

\newpage
\section*{Exercice 6 -- Exercice long type bac : étude d'un produit odorant \pts{12}}

Une espèce organique E est utilisée pour son odeur fruitée. Un élève cherche à l'identifier à partir de plusieurs informations expérimentales.

\subsection*{Document 1 -- Données sur l'espèce E}
\begin{itemize}[leftmargin=1.2em]
    \item La molécule E contient exactement \textbf{quatre atomes de carbone}.
    \item Son spectre IR présente une \textbf{bande intense vers \SI{1740}{cm^{-1}}}.
    \item Le spectre \textbf{ne présente pas} de large bande entre \SI{3200}{cm^{-1}} et \SI{3700}{cm^{-1}}.
    \item E a une odeur agréable et se dissout modérément dans l'eau.
\end{itemize}

\subsection*{Document 2 -- Espèces candidates}
\begin{enumerate}[label=\textbf{\Alph*.},leftmargin=1.2em]
    \item butan-1-ol
    \item butanal
    \item butanone
    \item éthanoate d'éthyle
\end{enumerate}

\subsection*{Document 3 -- Indications de cours}
\begin{itemize}[leftmargin=1.2em]
    \item Un alcool présente souvent une large bande O--H.
    \item Les aldéhydes, les cétones et les esters possèdent un groupe carbonyle.
    \item Les esters sont souvent associés à des odeurs fruitées.
\end{itemize}

\begin{enumerate}[label=\textbf{\arabic*.},leftmargin=1.2em]
    \item Parmi les quatre espèces candidates, lesquelles possèdent un groupe carbonyle ? \pts{1}
    \item Quelle espèce peut être éliminée immédiatement grâce à l'absence de large bande O--H ? \pts{1}
    \item Pourquoi la présence d'une bande vers \SI{1740}{cm^{-1}} ne permet-elle pas, à elle seule, de conclure que E est une cétone ? \pts{1.5}
    \item En exploitant les informations sur l'odeur, quelle famille est particulièrement suggérée ? \pts{1}
    \item Donner le nom et la famille de l'espèce D du document 2. \pts{1}
    \item Expliquer pourquoi l'espèce D est cohérente avec toutes les informations du document 1. \pts{2}
    \item Montrer que deux autres candidates peuvent sembler partiellement compatibles avec le spectre seul, mais pas avec l'ensemble des données. \pts{2}
    \item Conclure sur l'identité la plus probable de E. La conclusion est-elle absolument certaine ou seulement très probable ? Justifier. \pts{1.5}
    \item Le caractère ``modérément soluble dans l'eau'' est-il cohérent avec l'espèce retenue ? Expliquer qualitativement. \pts{1}
\end{enumerate}

\section*{Exercice 7 -- Question de synthèse et de réflexion \pts{6}}

Répondre en rédigeant soigneusement.

\begin{enumerate}[label=\textbf{\arabic*.},leftmargin=1.2em]
    \item Expliquer pourquoi \textbf{l'analyse d'un spectre IR seule} ne permet pas toujours d'identifier \emph{de manière unique} une molécule organique. \pts{2}
    \item Expliquer en quoi la \textbf{nomenclature} peut aider à interpréter ou vérifier les résultats d'une analyse spectrale. \pts{2}
    \item Donner un exemple simple montrant qu'une molécule peut contenir des liaisons polarisées tout en étant globalement apolaire. \pts{2}
\end{enumerate}

\newpage
\begin{center}
{\LARGE \textbf{Corrigé détaillé}}
\end{center}

\section*{Corrigé de l'exercice 1}
\begin{enumerate}[label=\textbf{\arabic*.},leftmargin=1.2em]
    \item Une large bande entre \SI{3200}{cm^{-1}} et \SI{3700}{cm^{-1}} est caractéristique d'une liaison O--H d'alcool. Son absence montre que le produit B ne possède pas de groupe hydroxyle d'alcool.

    \item La bande vers \SI{1715}{cm^{-1}} correspond à une liaison C=O. Il s'agit donc d'un \textbf{groupe carbonyle}.

    \item Un alcool doit présenter une bande O--H caractéristique. Comme cette bande est absente, B n'est plus un alcool. Cela est cohérent avec une transformation d'oxydation de l'alcool A vers une autre famille.

    \item Non. La présence d'un groupe carbonyle indique seulement que B appartient à une famille contenant C=O. Cela peut être un \textbf{aldéhyde} ou une \textbf{cétone} dans le cadre du programme étudié ici.

    \item Les deux espèces proposées, propanal et propanone, possèdent toutes les deux un groupe carbonyle et ne possèdent pas de groupe O--H. Le spectre donné est donc compatible avec \textbf{les deux}.

    \item Une information supplémentaire pourrait être la \textbf{formule semi-développée}, la \textbf{position du groupe carbonyle dans la chaîne}, une \textbf{réaction chimique caractéristique}, ou encore une autre donnée spectrale. En première, il suffit de dire qu'il faut une information permettant de savoir si le groupe carbonyle est en bout de chaîne (aldéhyde) ou au milieu (cétone).
\end{enumerate}

\section*{Corrigé de l'exercice 2}
\begin{enumerate}[label=\textbf{\arabic*.},leftmargin=1.2em]
    \item Une liaison polarisée est une liaison covalente dans laquelle les électrons liants sont attirés davantage par l'un des deux atomes. On obtient alors des charges partielles $\delta^+$ et $\delta^-$.

    \item Dans H$_2$O, les liaisons O--H sont polarisées car l'oxygène attire plus fortement les électrons. De plus, la molécule est coudée. Les effets des deux polarisations ne se compensent donc pas. La molécule possède un dipôle global non nul : elle est \textbf{polaire}.

    \item Dans CO$_2$, chaque liaison C=O est polarisée. Mais la molécule est \textbf{linéaire}. Les deux polarisations sont opposées et de même intensité ; elles se compensent. Ainsi, malgré des liaisons polarisées, la molécule peut être \textbf{globalement apolaire}.

    \item H$_2$S est plutôt \textbf{polaire}. En effet, les liaisons S--H sont polarisées et la molécule est coudée ; il n'y a donc pas compensation complète des dipôles de liaison.

    \item L'espèce qui interagit le plus fortement avec l'eau liquide est \textbf{H$_2$O elle-même}, puisqu'elle est très polaire et peut former de nombreuses liaisons hydrogène avec les autres molécules d'eau. Parmi les deux autres, H$_2$S interagit mieux avec l'eau que CO$_2$, car H$_2$S est polaire tandis que CO$_2$ est apolaire.
\end{enumerate}

\section*{Corrigé de l'exercice 3}
\begin{enumerate}[label=\textbf{\arabic*.},leftmargin=1.2em]
    \item 
    \begin{itemize}[leftmargin=1.2em]
        \item A : \textbf{éthanol}
        \item B : \textbf{propanal}
        \item C : \textbf{propanone}
        \item D : \textbf{acide propanoïque}
        \item E : \textbf{éthanoate d'éthyle}
    \end{itemize}

    \item 
    \begin{itemize}[leftmargin=1.2em]
        \item A : alcool
        \item B : aldéhyde
        \item C : cétone
        \item D : acide carboxylique
        \item E : ester
    \end{itemize}

    \item Les espèces possédant un groupe carbonyle sont \textbf{B, C, D et E}. En effet, aldéhydes, cétones, acides carboxyliques et esters contiennent tous la liaison C=O.
\end{enumerate}

\section*{Corrigé de l'exercice 4}
\begin{enumerate}[label=\textbf{\arabic*.},leftmargin=1.2em]
    \item La bande vers \SI{1720}{cm^{-1}} met en évidence un \textbf{groupe carbonyle C=O}.

    \item On élimine immédiatement la famille des \textbf{alcools}, car un alcool présenterait une bande large de liaison O--H.

    \item La formule brute C$_3$H$_6$O peut convenir à une molécule comportant un groupe carbonyle. Le programme de première permet d'envisager ici un \textbf{aldéhyde} ou une \textbf{cétone} à trois atomes de carbone.

    \item Les deux espèces encore possibles sont \textbf{le propanal} et \textbf{la propanone}.

    \item Non, le spectre donné ne suffit pas ici à les départager car ces deux espèces possèdent une liaison C=O et aucune liaison O--H. Le spectre permet d'identifier la famille large ``composé carbonylé'', mais pas toujours la sous-famille exacte.

    \item Le suffixe \textit{-one} désigne une \textbf{cétone}. Donc parmi les noms proposés, \textbf{propanone} est la cétone.

    \item On peut affirmer avec certitude que X n'est pas un alcool et qu'elle possède un groupe carbonyle. On peut dire qu'elle est compatible avec un aldéhyde ou une cétone, mais on ne peut pas identifier de manière unique la molécule avec les seules informations fournies.
\end{enumerate}

\section*{Corrigé de l'exercice 5}
\begin{enumerate}[label=\textbf{\arabic*.},leftmargin=1.2em]
    \item L'espèce la moins soluble dans l'eau paraît être \textbf{le butane}. C'est un hydrocarbure ne contenant pas d'atome très électronégatif, donc une molécule globalement apolaire, qui interagit peu avec l'eau, solvant polaire.

    \item Les espèces pouvant former des liaisons hydrogène avec l'eau sont surtout \textbf{butan-1-ol} et \textbf{acide éthanoïque}, car elles possèdent une liaison O--H. La \textbf{propanone} peut aussi interagir fortement avec l'eau grâce à l'oxygène du groupe carbonyle, qui peut participer à des interactions de type liaison hydrogène avec les molécules d'eau, même si la propanone ne possède pas elle-même de liaison O--H.

    \item La présence d'un groupe O--H favorise la solubilité dans l'eau car ce groupe est très polaire et permet l'établissement de liaisons hydrogène avec les molécules d'eau. Ces interactions stabilisent la dissolution.

    \item Oui. La propanone possède un oxygène dans son groupe carbonyle C=O. Ce groupe est polaire et l'oxygène peut interagir avec l'eau. La propanone peut donc être soluble dans l'eau même sans liaison O--H propre.

    \item Entre butan-1-ol et acide éthanoïque, l'\textbf{acide éthanoïque} semble interagir le plus fortement avec l'eau. Il possède à la fois un groupe carbonyle et une liaison O--H au sein du groupe carboxyle, ce qui rend l'espèce très polaire et particulièrement apte à établir des interactions fortes avec l'eau.
\end{enumerate}

\section*{Corrigé de l'exercice 6}
\begin{enumerate}[label=\textbf{\arabic*.},leftmargin=1.2em]
    \item Les espèces possédant un groupe carbonyle sont \textbf{B, C et D}, c'est-à-dire \textbf{butanal}, \textbf{butanone} et \textbf{éthanoate d'éthyle}.

    \item On élimine immédiatement \textbf{A : butan-1-ol}, car un alcool présenterait une large bande O--H entre environ \SI{3200}{cm^{-1}} et \SI{3700}{cm^{-1}}.

    \item La bande vers \SI{1740}{cm^{-1}} révèle la présence d'un groupe carbonyle, mais ce groupe n'est pas réservé aux cétones. Les aldéhydes et les esters possèdent eux aussi une liaison C=O. On ne peut donc pas conclure ``cétone'' à partir de cette seule bande.

    \item Les odeurs fruitées suggèrent particulièrement la famille des \textbf{esters}.

    \item L'espèce D est \textbf{l'éthanoate d'éthyle}. Sa famille est celle des \textbf{esters}.

    \item L'éthanoate d'éthyle est cohérent avec toutes les données :
    \begin{itemize}[leftmargin=1.2em]
        \item il contient quatre atomes de carbone ;
        \item il possède un groupe carbonyle, donc une bande intense vers \SI{1740}{cm^{-1}} est cohérente ;
        \item il ne possède pas de groupe O--H d'alcool, donc l'absence de large bande O--H convient ;
        \item les esters sont souvent associés à des odeurs fruitées ;
        \item une solubilité modérée dans l'eau est plausible pour une petite molécule ester : il existe des interactions possibles avec l'eau, mais elles sont moins fortes que pour un alcool ou un acide carboxylique.
    \end{itemize}

    \item Le \textbf{butanal} et la \textbf{butanone} peuvent sembler compatibles avec le spectre seul, car ce sont aussi des composés carbonylés sans liaison O--H. En revanche, l'information sur l'odeur fruitée oriente nettement vers un ester. De plus, la cohésion avec toutes les données expérimentales est meilleure pour l'éthanoate d'éthyle.

    \item L'identité la plus probable de E est \textbf{l'éthanoate d'éthyle}. La conclusion est \textbf{très probable} plutôt qu'absolument certaine, car l'identification complète d'une molécule demanderait en général davantage d'informations expérimentales. Mais parmi les candidates proposées, l'ester est clairement la meilleure réponse.

    \item Oui, cela est cohérent. Un ester est une molécule polaire grâce à la présence du groupe carbonyle et de l'atome d'oxygène, mais il ne possède pas de groupe O--H. Ses interactions avec l'eau existent donc, sans être aussi fortes que celles d'un alcool ou d'un acide. Une \textbf{solubilité modérée} est donc plausible.
\end{enumerate}

\section*{Corrigé de l'exercice 7}
\begin{enumerate}[label=\textbf{\arabic*.},leftmargin=1.2em]
    \item Un spectre IR renseigne surtout sur la présence de certains \textbf{groupes caractéristiques}. Or plusieurs molécules différentes peuvent contenir le même groupe. Par exemple, un aldéhyde et une cétone possèdent tous deux une liaison C=O. On peut donc souvent identifier une \textbf{famille} ou une \textbf{partie de structure}, mais pas toujours la molécule unique.

    \item La nomenclature donne des indices directs sur la famille chimique : \textit{-ol}, \textit{-al}, \textit{-one}, \textit{-oïque}, etc. Elle permet donc de vérifier si les bandes observées sont cohérentes avec le nom proposé. Par exemple, une espèce nommée \textit{propanone} doit appartenir à la famille des cétones et donc posséder un groupe carbonyle, sans groupe O--H d'alcool.

    \item Un exemple simple est le \textbf{dioxyde de carbone CO$_2$}. Les liaisons C=O sont polarisées, mais la molécule est linéaire, donc les dipôles de liaison se compensent. La molécule est alors globalement \textbf{apolaire} malgré des liaisons polarisées.
\end{enumerate}

\vfill

\end{document}
